\chapter{Exchanges, Brokers, Vendors and Regulators as Employers}\label{ch:in:exchanges-brokers-vendors-and-regulators-as-employers}

Listed exchange groups publish what private trading firms do not: their revenue, their operating income and the number
of people they employ, every year. In 2025 the four largest US groups employed between 1\,661 and 12\,844 people each and
earned between \$0.24 million and \$1.09 million of operating income per employee. The market's regulator publishes
even more: the \omterm{def:in:exchanges-brokers-vendors-and-regulators-as-employers:band}{pay band} of every grade of its staff and the supplement for each city, up to a cap of \$292\,300 in 2026.
The rest of the ecosystem --- exchanges, clearing houses, brokers, data and technology vendors, regulators --- is the most
transparent part of the industry to work in, and the one whose jobs are closest to engineering and public service. This
chapter describes it as an employer.

\section{Exchange groups and clearing houses}

\begin{definition}[Financial market infrastructure]\label{def:in:exchanges-brokers-vendors-and-regulators-as-employers:fmi}
A \emph{financial market infrastructure}\index{financial market infrastructure} is a multilateral system among its
participants for clearing, settling or recording payments, securities, derivatives or other transactions: a
systemically important payment system, a central securities depository, a securities settlement system, a central
counterparty or a trade repository. Exchanges and trading venues are the markets these systems serve and are often owned
by the same groups.
\end{definition}

An exchange group sells trading, clearing, data, index licences and connectivity (Book~1, chapter~4; Book~14). It is a
technology company regulated as a utility: its engineers build and run matching engines, market-data systems and
clearing platforms (Book~10, chapter~27, builds a simulator of one); its quants design margin models, index methodologies
and surveillance; its product staff design contracts and fee schedules. The work is on the other side of the trading
firms' problems: the latency they race for is the latency an exchange engineer measures and makes fair.

\begin{dated}{2025-12}{Four exchange groups in their filings}\label{dat:in:exchanges-brokers-vendors-and-regulators-as-employers:exchanges}
\footnotesize
\begin{tabular}{@{}lrrrrr@{}}
\toprule
group (10-K 2025) & revenue, \$m & op.\ income, \$m & staff & revenue/head & op.\ income/head\\
\midrule
CME Group & 6\,520.6 & 4\,229.5 & about 3\,875 & \$1.68m & \$1.09m\\
Cboe Global Markets & 4\,714.2 & 1\,467.1 & 1\,661 & \$2.84m & \$0.88m\\
Intercontinental Exchange & 12\,640 & 4\,929 & 12\,844 & \$0.98m & \$0.38m\\
Nasdaq & 8\,262 & 2\,331 & 9\,525 & \$0.87m & \$0.24m\\
\bottomrule
\end{tabular}

\smallskip\noindent Revenue is each group's total revenue as reported, which for three of them includes transaction-based
expenses (rebates and fees passed through to participants); operating income is after them. Two of the groups report segments outside
trading: one reports exchanges, fixed income and data services, and mortgage technology; another capital access
platforms, financial technology and market services.
\end{dated}

The four differ more in operating income per head than in anything else
(\cref{fig:in:exchanges-brokers-vendors-and-regulators-as-employers:perhead}): the group whose revenue is mostly derivatives
trading and clearing earns \$1.09 million per employee, 4.5 times
the group that reports a financial-technology segment.
Revenue per head is less comparable, because pass-through fees inflate the revenue of groups that pay rebates. For an
employee the difference is the business: a derivatives exchange employs people close to trading and clearing, a
diversified group many more in software, data and services.

\begin{omfigure}
\begin{tikzpicture}
\begin{axis}[width=10.5cm,height=5.2cm,xbar,area legend,bar width=7pt,xmin=0,xmax=3.2,ymin=-0.6,ymax=3.6,xlabel={\small \$ million per employee, 2025},
  ytick={0,1,2,3},yticklabels from table={figdata/industry/09-exchanges-brokers-vendors-and-regulators-as-employers/perhead.csv}{firm},
  y dir=reverse,tick label style={font=\footnotesize},grid=major,grid style={gray!20},table/col sep=comma,
  legend style={legend cell align=left,font=\scriptsize,at={(0.5,-0.3)},anchor=north,/tikz/every even column/.append style={column sep=8pt}},legend columns=2]
\addplot[fill=omDef!55,draw=omDef] table[y=k,x=opincome]{figdata/industry/09-exchanges-brokers-vendors-and-regulators-as-employers/perhead.csv};
\addplot[fill=gray!40,draw=gray] table[y=k,x=revenue]{figdata/industry/09-exchanges-brokers-vendors-and-regulators-as-employers/perhead.csv};
\legend{operating income per head,revenue per head}
\end{axis}
\end{tikzpicture}

\omcaption{Operating income and revenue per employee of four US exchange groups, 2025, ordered by operating income per
head. Revenue includes transaction-based expenses for three of the groups. Data: Forms 10-K for 2025 and SEC XBRL
company facts, through \texttt{in\_infra.exchanges}.}\label{fig:in:exchanges-brokers-vendors-and-regulators-as-employers:perhead}
\end{omfigure}

\section{Brokers and trading-technology vendors}

Brokers sell access and execution: agency brokers route and work orders (Book~10, chapter~20), prime brokers finance and
custody (Book~1, chapter~6), clearing brokers stand between their clients and the clearing houses (Book~1, chapter~30).
Trading-technology vendors sell the systems that firms do not build themselves: order and execution management
systems, market-data feeds and ticker plants, risk and compliance software, connectivity (Book~14, chapter~27, maps
them). For a quant or an engineer these employers offer the same technical problems as a trading firm --- latency,
correctness, scale --- with the firm's customers' P\&L instead of its own at stake, and pay that follows a software
company's economics: salaries and equity rather than a share of trading profit.

\section{Data vendors}

Market-data vendors collect, normalise, store and redistribute data: exchange feeds, reference data, prices of
instruments that trade over the counter, alternative data (Book~7, chapter~12). Their quantitative jobs are in data
quality, evaluated pricing (models that price bonds and derivatives that did not trade), index construction and
analytics. Their engineering jobs are those of a large-scale data platform (Book~15). The work is closest to what a
research platform team does inside a trading firm, but for many clients at once.

\section{Regulators and central banks}

Regulators employ economists, quantitative analysts, examiners, lawyers and technologists. Their market-structure and
surveillance teams use the methods of this series on the market as a whole: the race study of chapter~2 is regulators'
research, and surveillance uses the detectors of Book~9 (chapter~29) on the consolidated audit trail. Central banks
employ quants in market operations, reserve management and financial-stability analysis. The pay is public.

\begin{definition}[Pay band]\label{def:in:exchanges-brokers-vendors-and-regulators-as-employers:band}
A \emph{pay band}\index{pay band} is the published range of basic pay for a grade of an employer's pay scale, from a
minimum to a maximum; a public employer's scale also publishes the supplements (for example, by city) that apply to it
and any cap on total pay.
\end{definition}

\begin{dated}{2026-04}{One regulator's pay scale}\label{dat:in:exchanges-brokers-vendors-and-regulators-as-employers:sec}
The SEC's 2026 basic pay for its SK scale, effective 5 April 2026, runs from \$24\,728--\$36\,213 at grade 1 to
\$138\,383--\$239\,449 at grade 17 (supervisory). Selected grades: SK-12 \$83\,718--\$141\,843; SK-13 \$99\,557--\$168\,668;
SK-14 \$114\,315--\$193\,670; SK-16 \$130\,550--\$221\,711. Locality rates apply to basic pay by duty station: New York
37.95\%, San Francisco 46.34\%, Washington DC 33.94\%, Chicago 30.86\%. ``The salary cap, including locality pay, is
\$292\,300 for all SK staff.''
\end{dated}

\begin{method}[A public band as an annual range]\label{met:in:exchanges-brokers-vendors-and-regulators-as-employers:band}
\begin{enumerate}
\item Take the grade's minimum and maximum of basic pay.
\item Multiply both by one plus the duty station's locality rate.
\item Cap each end at the scale's cap.
\item Compare the resulting range with the range another source gives for the same occupation (chapter~14), by the
share of the other range the band covers.
\end{enumerate}
\end{method}

In New York, an SK-14 position pays \$157\,698 to \$267\,168; an SK-16 position \$180\,094 to \$292\,300, the upper end
capped (\cref{fig:in:exchanges-brokers-vendors-and-regulators-as-employers:bands}). The bands are wide --- the top of a
grade is about 1.7 times its bottom --- because they cover a career's progression within a grade, and a new hire is
usually placed near the bottom. The cap binds the top grades: from SK-16 up, the upper end in New York is the cap, not
the grade.

\begin{omfigure}
\begin{tikzpicture}
\begin{axis}[width=10.5cm,height=5.6cm,xmin=8.4,xmax=17.6,ymin=0,ymax=340,xtick={9,10,11,12,13,14,15,16,17},
  xlabel={\small SK grade},ylabel={\small annual pay in New York (\$ thousand)},tick label style={font=\footnotesize},
  grid=major,grid style={gray!20},table/col sep=comma]
\addplot[only marks,mark=-,omDef,thick,error bars/.cd,y dir=both,y explicit,error bar style={line width=5pt,omDef!45}]
  table[x=grade,y=mid,y error=half]{figdata/industry/09-exchanges-brokers-vendors-and-regulators-as-employers/skbands.csv};
\draw[omThm,dashed] (axis cs:8.4,292.3) -- (axis cs:17.6,292.3);
\node[font=\scriptsize,omThm,anchor=south west] at (axis cs:8.5,292.3) {cap \$292\,300};
\end{axis}
\end{tikzpicture}

\omcaption{The SEC's 2026 \omterm{def:in:exchanges-brokers-vendors-and-regulators-as-employers:band}{pay bands} for grades SK-9 to SK-17 in New York: basic pay times one plus the 37.95\% locality
rate, capped at \$292\,300 (dashed). Each bar runs from the grade's minimum to its maximum. Data: SEC compensation overview,
through \texttt{firm.payband.at}.}\label{fig:in:exchanges-brokers-vendors-and-regulators-as-employers:bands}
\end{omfigure}

\section{What the infrastructure offers in the filings}

The four exchange groups also file labour condition applications (chapter~14), and their offers can be set beside the
banks' and the market makers', and beside the regulator's band.

\begin{dated}{2025-09}{Exchange groups in the filings}\label{dat:in:exchanges-brokers-vendors-and-regulators-as-employers:lca}
Fiscal 2025 labour condition applications, median offered base: exchange groups, all titles \$130\,800 (462 applications,
4 employers), software \$140\,400 (329), quantitative research \$107\,100 (35), machine learning and data \$108\,500 (60);
banks \$154\,260, \$155\,688, \$158\,100 and \$140\,714; market makers \$175\,000 in each of the four. Base salary only.
\end{dated}

\begin{omfigure}
\begin{tikzpicture}
\begin{axis}[ybar,area legend,width=10.5cm,height=5.4cm,xmin=0.5,xmax=4.5,ymin=0,ymax=200,bar width=10pt,xtick={1,2,3,4},
  xticklabels from table={figdata/industry/09-exchanges-brokers-vendors-and-regulators-as-employers/filings.csv}{role},
  ylabel={\small median offered base, \$ thousand},tick label style={font=\footnotesize},grid=major,grid style={gray!20},
  table/col sep=comma,legend style={legend cell align=left,font=\scriptsize,at={(0.5,-0.25)},anchor=north,legend columns=3,
  /tikz/every even column/.append style={column sep=8pt}}]
\addplot[fill=black!30,draw=none] table[x=pos,y=exchange]{figdata/industry/09-exchanges-brokers-vendors-and-regulators-as-employers/filings.csv};
\addplot[fill=omDef,draw=none] table[x=pos,y=bank]{figdata/industry/09-exchanges-brokers-vendors-and-regulators-as-employers/filings.csv};
\addplot[fill=omThm!70,draw=none] table[x=pos,y=market_maker]{figdata/industry/09-exchanges-brokers-vendors-and-regulators-as-employers/filings.csv};
\draw[black!60,dashed] (axis cs:0.5,157.698) -- (axis cs:4.5,157.698);
\legend{exchange groups,banks,market makers}
\end{axis}
\end{tikzpicture}

\omcaption{Median offered base in fiscal 2025 labour condition applications by role family and kind of employer, with
the bottom of the SEC's SK-14 band in New York (dashed). Data: \texttt{data/industry/lca\_ranges.csv}, through
\texttt{in\_infra.filings}.}\label{fig:in:exchanges-brokers-vendors-and-regulators-as-employers:filings}
\end{omfigure}

The exchanges' offers sit below the banks' in every family, and well below the market makers'
(\cref{fig:in:exchanges-brokers-vendors-and-regulators-as-employers:filings}): \$130\,800 against \$154\,260 and
\$175\,000 for all titles; for research, \$107\,100 against \$158\,100 and \$175\,000. The regulator's band sits between:
an SK-14 position in New York starts at \$157\,698 and runs to \$267\,168, above the exchanges' median offers and in
the range of the banks'. None of these counts bonus, which is larger at the trading firms (chapter~14); the public band
is the whole of the regulator's pay.

\section{Moving between the sides}

People move between the industry and its infrastructure in both directions: engineers from exchanges to trading firms
and back, quants from regulators to banks and from banks to regulators, market-structure economists from venues to
regulators. Moves out of a regulator are subject to the post-employment rules of public service, and moves into one to
conflict-of-interest rules; both are the employer's rules, set out in its own codes. For a career, the infrastructure is
where one learns how the market works from the inside --- its plumbing, its rules and its data --- at pay that is
published and more stable, and usually lower at the top than the trading firms' (chapter~14).

\section{Tutorial: revenue per head at the infrastructure}\label{tut:in:exchanges-brokers-vendors-and-regulators-as-employers:tut}

\begin{tutorial}
\textbf{Goal.} Compute exchange groups' revenue and operating income per head from their filings, and turn a regulator's
published scale into annual ranges. \textbf{End state:}
\cref{fig:in:exchanges-brokers-vendors-and-regulators-as-employers:perhead,fig:in:exchanges-brokers-vendors-and-regulators-as-employers:bands}.
\begin{tutsteps}
\item \textbf{The filings.} \texttt{data/industry/infrastructure\_employers.csv} holds each group's total revenue and
operating income (from SEC XBRL company facts where tagged, else the income statement) and its year-end employees (from
the 10-K text); \texttt{in\_infra.exchanges()} computes the per-head measures and the operating margin.
\item \textbf{The scale.} \texttt{firm.payband.load\_scale} reads the SK grades, the locality rates and the cap;
\texttt{firm.payband.at(scale, grade, station)} returns the annual range
(\cref{lst:in:exchanges-brokers-vendors-and-regulators-as-employers:at}).
\omcode{code/firm/payband/firm_payband.py}{41}{52}{A grade's annual range at a duty station, and the overlap of two
ranges.}\label{lst:in:exchanges-brokers-vendors-and-regulators-as-employers:at}
\item \textbf{Compare.} \texttt{firm.payband.overlap(a, b)} gives the share of a range $b$ that a band $a$ covers, for
chapter~14's comparison with labour-condition filings.
\end{tutsteps}
\end{tutorial}

\textbf{What to change next.} Compute each group's revenue per head net of transaction-based expenses where the 10-K
reports it (exercise~7); tabulate the SK bands for Chicago and Washington and compare.

\section{Build: public pay scales}\label{bld:in:exchanges-brokers-vendors-and-regulators-as-employers:payband}

\begin{build}
\bfield{Purpose} Read published public-sector pay scales as annual ranges, so that chapter~14 can set them beside the
private-sector ranges of labour-condition filings and chapter~24 beside control functions' pay.
\bfield{Interface} \texttt{firm.payband}: \texttt{Band(grade, lo, hi)}, \texttt{Scale(name, bands, locality, cap)};
\texttt{load\_scale(name, bands\_csv, locality\_csv, cap)}; \texttt{at(scale, grade, station)}; \texttt{overlap(a, b)}.
\bfield{Rules} Locality applies to basic pay; the cap applies to the total, at each end of the range; an unknown station
or grade is an error.
\bfield{Acceptance tests} \texttt{code/firm/payband/tests/}: locality and cap on a constructed scale; overlap of
disjoint and overlapping ranges; loading from CSV.
\bfield{Stretch} Step tables within a grade (pay by years of service); several scales in one currency for cross-country
comparison (chapter~27).
\end{build}

\begin{omsources}
\item Forms 10-K for 2025 of CME Group, Intercontinental Exchange, Nasdaq and Cboe Global Markets; SEC XBRL company facts.
\item SEC, Compensation overview (2026 SK pay scale and locality rates).
\item CPMI-IOSCO, Principles for financial market infrastructures (2012), BIS publication page.
\end{omsources}

\section{Exercises}

\begin{exercise}[$\star$]\label{exo:in:exchanges-brokers-vendors-and-regulators-as-employers:1}
Compute each exchange group's operating margin in 2025.
\end{exercise}

\begin{exercise}[$\star$]\label{exo:in:exchanges-brokers-vendors-and-regulators-as-employers:2}
Compute the New York range of an SK-13 position, and its ratio of top to bottom.
\end{exercise}

\begin{exercise}[$\star$]\label{exo:in:exchanges-brokers-vendors-and-regulators-as-employers:3}
Which grades' New York upper end is set by the cap rather than the grade?
\end{exercise}

\begin{exercise}[$\star\star$]\label{exo:in:exchanges-brokers-vendors-and-regulators-as-employers:4}
Why is revenue per head a poorer comparison across the four exchange groups than operating income per head?
\end{exercise}

\begin{exercise}[$\star\star$]\label{exo:in:exchanges-brokers-vendors-and-regulators-as-employers:5}
An SK-14 position in Chicago and one in New York: compute both ranges and the difference in their midpoints.
\end{exercise}

\begin{exercise}[$\star\star$]\label{exo:in:exchanges-brokers-vendors-and-regulators-as-employers:6}
From \cref{fig:in:exchanges-brokers-vendors-and-regulators-as-employers:perhead}, which group has the widest gap
between revenue and operating income per head, and what might explain it?
\end{exercise}

\begin{exercise}[$\star\star\star$]\label{exo:in:exchanges-brokers-vendors-and-regulators-as-employers:7}
\emph{Coding.} Nasdaq reports total revenues less transaction-based expenses of \$5\,249 million. Recompute its revenue per
head on that basis and its operating margin on it. How much of the gap in operating margin with the derivatives exchange closes?
\end{exercise}

\begin{exercise}[$\star\star\star$]\label{exo:in:exchanges-brokers-vendors-and-regulators-as-employers:8}
\emph{Find the flaw.} ``An SK-16 examiner in New York earns \$292\,300 and a quant at a trading firm earns less than that
in base salary, so the regulator pays more.''
\end{exercise}

\section{Problem: Revenue per Head at the Infrastructure}

\begin{problem}[{Weekend problem --- revenue per head at the infrastructure}]\label{pb:in:exchanges-brokers-vendors-and-regulators-as-employers:1}
An engineer with offers from an exchange group and a regulator wants to understand both employers from their public
numbers.

\medskip\noindent\textbf{Part I --- The infrastructure.}
\begin{enumerate}
\item Define a \omterm{def:in:exchanges-brokers-vendors-and-regulators-as-employers:fmi}{financial market infrastructure} and list its five kinds.
\item What does an exchange group sell, and which quant and engineering jobs does it employ?
\item Give the four groups' employees in 2025.
\item Give their revenue and operating income per head.
\item What inflates the revenue of three of them, and how?
\end{enumerate}

\medskip\noindent\textbf{Part II --- The comparison.}
\begin{enumerate}[resume]
\item Give the four operating margins.
\item How many times the lowest is the highest operating income per head?
\item Which segments outside trading do two of the groups report?
\item What does a broker sell, and a data vendor?
\item What distinguishes a trading-technology vendor's pay from a trading firm's?
\end{enumerate}

\medskip\noindent\textbf{Part III --- The regulator.}
\begin{enumerate}[resume]
\item Define a \omterm{def:in:exchanges-brokers-vendors-and-regulators-as-employers:band}{pay band}.
\item State the method that turns a band into an annual range.
\item Give the New York ranges of SK-12, SK-14 and SK-16.
\item Which grades are capped in New York, and at what?
\item Why are the bands so wide, and where is a new hire placed?
\end{enumerate}

\medskip\noindent\textbf{Part IV --- The verdict.}
\begin{enumerate}[resume]
\item State the \emph{named result}: the range of operating income per head across the four groups, and the New York
range of an SK-14 position.
\item What rules govern moves out of and into a regulator?
\item What does the infrastructure teach that a trading firm does not?
\item Where does chapter~14 compare these bands with private pay?
\item In two sentences, how should the engineer compare the two offers?
\end{enumerate}
\end{problem}

\section{Interview questions}

\begin{interviewq}[$\star$ \iqroles{developer}]\label{iq:in:exchanges-brokers-vendors-and-regulators-as-employers:1}
What does an exchange's matching-engine team worry about that a trading firm's strategy engineers do not?
\end{interviewq}

\begin{interviewq}[$\star$ \iqroles{risk}]\label{iq:in:exchanges-brokers-vendors-and-regulators-as-employers:2}
What does a clearing house's margin model have to do, and who bears the risk if it is wrong?
\end{interviewq}

\begin{interviewq}[$\star\star$ \iqroles{developer}]\label{iq:in:exchanges-brokers-vendors-and-regulators-as-employers:3}
An exchange wants every participant in a colocation hall to receive market data at the same time. How would you design
and test that?
\end{interviewq}

\begin{interviewq}[$\star\star$ \iqroles{researcher}]\label{iq:in:exchanges-brokers-vendors-and-regulators-as-employers:4}
A regulator asks you to estimate how much of a market's volume is in latency races. What data do you need, and what is
the first thing you measure?
\end{interviewq}

\begin{interviewq}[$\star\star$ \iqroles{researcher}]\label{iq:in:exchanges-brokers-vendors-and-regulators-as-employers:5}
A data vendor prices a corporate bond that has not traded for three weeks. How, and how do you know the price is good?
\end{interviewq}

\begin{interviewq}[$\star\star\star$ \iqroles{risk, developer}]\label{iq:in:exchanges-brokers-vendors-and-regulators-as-employers:6}
An exchange pays rebates to liquidity providers and reports them as transaction-based expenses. How do rebates change the
meaning of its revenue, and how would you compare two exchanges fairly?
\end{interviewq}
